\subsection{The E6 infinity-coordinate disc and its full marked comparison}
\begin{statement}{P3K.97}{Proved}
Let \(w\) vary on a sufficiently small disc about either point \(w^2=3\), and set
\[
 D=w^2-3,\quad R=3-2w^2,\quad v=w^2-1,\quad
 r=R/D^2,\quad t_c=v/D^3.
\]
On its punctured disc retain the actual original source
\[
 y^2=x^3+r x^2+t x+t^2,\qquad G=(0,t),\qquad
 M_a=\mathcal O(aG-O),\qquad Q=bG,
\]
with the whole selected zero \(t=t_c\), the full Cartier suspension
\(B_{\rm sel}=(t-t_c)B_{a,b}/t^3\), and either complete cap progression
\[
 (a,b;m)=(6,1;3l+2),\ l\ge0,\qquad
 (a,b;m)=(2,3;3l),\ l\ge1.
\]
Use either signed primitive norm \(n=\pm1\), the prescribed finite
character \(\zeta_m^d\), a fixed sufficiently small nonzero scalar,
the entire reference annular map, and no additional logarithmic
sites. For every prescribed lattice-basis fan at the selected cusp,
this punctured family is isomorphic, by the actual source, Cartier,
root, physical-circle and annular maps of the proof, to the restriction
of a smooth proper oriented sphere family on the full native disc
\[
 Y^2=X^3+R X^2+DTX+T^2,\qquad G=(0,T),\qquad q=T-v.
\]
The comparison retains both point and tensor Tate frames, the full
finite-root character, every period and all parameter scalar factors.
It intertwines the specified forward endpoints and hence their inverse
physical closings. The actual positive marked \(w\)-loop has return
\(0\in\Theta_7\).

The original infinity parameter is exactly
\(\epsilon=1/r=D^2/(3-2w^2)\). Each such positive \(w\)-loop
has degree two over the original positive infinity meridian.
No shorter original-parameter endpoint or return is inferred.
\end{statement}
\begin{proof}
\par

\medskip\noindent\textbf{DZ6.1. The native disc and actual point-line maps}

\par\nobreak

Put
\[
 D=w^2-3,\qquad R=3-2w^2,\qquad v=w^2-1,\qquad
 s_0=D^2/4-R=(w^2-1)(w^2+3)/4.
\]
On a small disc about either of the points \(w^2=3\), use
\[
 Y^2=X^3+R X^2+DTX+T^2,\qquad G=(0,T),\qquad q=T-v.
 \tag{DZ6.1}
\]
Here \(s_0=3\), \(v=2\) at the centre. The selected nodal
fibre remains disjoint from the I2 fibre at \(T=0\). The
actual source A1 coordinates are
\[
 U=Y-T-DX/2,\quad V=Y+T+DX/2,\quad UV=X^2(X-s_0).
\]
Thus the same simultaneous A1 resolution, using the extending
unit square root of \(X-s_0\), is defined on the whole disc.
The exact Miller functions are
\[
 f_2=\frac{U}{X-s_0},\qquad
 f_3=\frac{s_0 XU}{s_0V-(Ds_0+2T)X},\qquad
 f_6=f_3^2 f_{2,3G}.
 \tag{DZ6.2}
\]
For \((a,b)=(6,1)\) retain \(B_6\); for \((a,b)=(2,3)\)
retain the full three-step product of EX.1. In both cases the
selected whole-zero coefficient is exactly
\[
                         B_{\rm sel}=\frac{q}{T^3}B_{a,b}.
 \tag{DZ6.3}
\]
The generic component valuations are unchanged because \(s_0\)
is a unit. The explicit source translations and line maps extend
on the native disc by those same actual resolution charts.

\par

\medskip\noindent\textbf{DZ6.2. All actual roots and their old/native coordinate arrows}

\par\nobreak

For good roots \(m=3l\), \(l\ge1\), the native root is
\[
 Y'^2=X'^3+Rz^{2l}X'^2+Dz^lX'+1.
\]
For reduced roots \(m=3l+2\), \(l\ge0\), it is
\[
 Y'^2=X'^3+Rz^{2l+2}X'^2+Dz^{l+2}X'+z^2.
 \tag{DZ6.4}
\]
The full EX middle-chart formula applies with
\(b_0=Dz^l,c_0=Rz^{2l}\); its exact universal polynomial
identities therefore retain these parameter functions. The marked
sections, division sections and all graph ideals use \(s_0\).
Both central selected root multipliers are four. The actual
starting line map is the same \(F=\tau g_{P,R_{\rm div}}\)
from EX.15, with its entire chord and coordinate factor.

On \(D\ne0\), the original source coordinates are related by
\[
 X=D^2x,\quad Y=D^3y,\quad T=D^3t,\quad
 r=R/D^2,\quad s=s_0/D^2,\quad q_{\infty,\rm old}
       =D^3q_{\infty,\rm new}.
 \tag{DZ6.5}
\]
Every original \(f_2\) has weight \(D^{-1}\), hence the
six-factor coefficient has weight \(D^{-6}\). The factor
\((t-t_c)/t^3\) has weight \(D^6\). Consequently the full
suspension scalar is unchanged by DZ6.5.

The selected toroidal arrow is also literal. On its nodal
normalization \(S_{\rm new}=D S_{\rm old}\) and
\(w_{\rm old}=w/D\), so
\[
 Z_{\rm old}=\frac{S_{\rm old}-w/D}{S_{\rm old}+w/D}
            =Z_{\rm new},\qquad q_{\rm old}=D^{-3}q_{\rm new}.
 \tag{DZ6.5a}
\]
Away from the node the same equality is obtained by integrating
the actual invariant differential divided by its transported
vanishing period; both scale by the same factor under DZ6.5.
Thus the full Tate period and marked point agree.
Let \(\sigma_a\) be the canonical meromorphic section of
\(\mathcal O(aG-O)\), and use the actual theta frame
\[
 e_\theta=\sigma_a/F_a,\qquad
 F_a(Z)=\Theta(Z/P_G^a)/\Theta(Z).
\]
The Cartier pullback sends \(\sigma_{a,\rm old}\) to
\(\sigma_{a,\rm new}\), while the two complete theta products
have identical arguments and periods. Hence this exact frame
map gives \(V_{\rm old}=V_{\rm new}\).
For example, a local Cartier equation
\(S_{\rm old}-S_{a,\rm old}=D^{-1}(S_{\rm new}-S_{a,\rm new})\)
changes its regular frame by \(D\); the two factors cancel in
the displayed \(e_\theta\). This is the actual frame arrow,
including its coefficient, rather than an omitted unit.
Consequently each toric character has the complete pullback
\[
 Z_{\rm old}^pV_{\rm old}^j q_{\rm old}^h
       =D^{-3h}Z_{\rm new}^pV_{\rm new}^j q_{\rm new}^h.
 \tag{DZ6.5b}
\]
The factor \(D^{-3h}\) and its parameter winding are retained.
On the parameter annulus it is an extending chart unit, so the
same fan and physical circle are carried by both inverse maps.

This point frame has the following exact dictionary with the
earlier point-to-tensor Tate frame. Write \(v_M\) for the
coefficient in the canonical meromorphic section \(\sigma_a\),
and set \(f_P=\Theta(Z/P_G)/\Theta(Z)\),
\(f_{P^a}=\Theta(Z/P_G^a)/\Theta(Z)\). Then
\[
 V_{\rm pt}=f_{P^a}(Z)v_M,\qquad
 V_{\rm tensor}=f_P(Z)^a f_a(p)^{-1}v_M,\qquad
 \frac{V_{\rm tensor}}{V_{\rm pt}}
       =\frac{f_P^a}{f_a f_{P^a}}.
 \tag{DZ6.5c}
\]
The last ratio is the full nonzero base constant: the numerator
and denominator have the same divisor and the same Tate
automorphy factor, so their quotient is a holomorphic unit on
the compact elliptic fibre. Its actual value, supplied by the
displayed functions, is retained. Since the Miller arrow has
the exact weight \(f_{a,\rm old}=D^{1-a}f_{a,\rm new}\),
the point coordinate is unchanged whereas
\[
                 V_{{\rm tensor},\rm old}
                      =D^{a-1}V_{{\rm tensor},\rm new}.
 \tag{DZ6.5d}
\]
Thus the complete Cox-character map in this tensor-reference
frame is
\[
 Z_{\rm old}^pV_{{\rm tensor},\rm old}^j q_{\rm old}^h
  =D^{(a-1)j-3h}
       Z_{\rm new}^pV_{{\rm tensor},\rm new}^j q_{\rm new}^h.
 \tag{DZ6.5e}
\]
The physical-line factors are \(D^5\) for \(M_6\) and \(D\)
for \(M_2\). They are closed nonzero parameter units on the
punctured comparison, with their full winding retained.

Write \(\eta=\log D/(2\pi i)\) on its parameter cover.
Let \(J_\eta\) denote the \textbf{new-to-old} root coordinate map.
For good roots it is
\[
 z_{\rm old}=D^{3/m}z_{\rm new},\quad
 X'_{\rm old}=X'_{\rm new},\quad Y'_{\rm old}=Y'_{\rm new}.
\]
For reduced roots it is
\[
 z_{\rm old}=D^{3/m}z_{\rm new},\quad
 X'_{\rm old}=D^{2/m}X'_{\rm new},\quad
 Y'_{\rm old}=D^{3/m}Y'_{\rm new},
\]
\[
 \tau'_{\rm old}=D^{-1/m}\tau'_{\rm new},\qquad
 \tau_{0,\rm old}=D^{-2}\tau_{0,\rm new}.
 \tag{DZ6.6}
\]
These maps are the actual base-change and minimal-coordinate maps,
not equalities of fibre types. In the reduced branch the root
\(B'_6\) has weight \(D^{6/m}\), canceled by
\(z_{\rm old}^{-2}\). The chord-over-vertical function has
weight \(D^{1/m}\), canceled by \(\tau'\). Thus
\(F\), its full norm ratio and its reference-normalized root
are literally carried to the corresponding native maps.
The good-root native point-line \(N_2\) coordinates already
agree. The common-graph correction still carries its full
monomial factor: \(u_{\rm old}=D u_{\rm new}\), so an
actual correction generator \(w_0^r u^s\) pulls back with
the coefficient \(D^s\). These factors remain on both
projections of the Cartier graph. They do not alter the
displayed point-line coordinate equality.

Let \(\mathcal D\) be the actual finite root deck. For the
reduced branch the identities \(6k\equiv-2\pmod m\),
\(9k\equiv-3\pmod m\), \(k=2l+1\), show that
\[
 J_{\eta+1}=\mathcal D_{\rm old}^{\,3}J_\eta,
 \qquad J_{\eta+1}^{-1}
       =\mathcal D_{\rm new}^{-3}J_\eta^{-1}.
 \tag{DZ6.7}
\]
The same typed equations hold for the good branch. The inverse
arrow therefore has the inverse deck; these are two directions
of the same coordinate square.

\par

\medskip\noindent\textbf{DZ6.3. Full cap endpoint, including scalar and finite character}

\par\nobreak

Keep \(R_\lambda^m=\lambda\), the actual signed norm \(n\),
finite character \(\zeta_m^d\), and the full EX norm lift
\(\widetilde F\). On the native root put
\[
 A=R_\lambda^n\zeta_m^d\widetilde F^{\,n}\mathcal D,
 \qquad B=\lambda B_{\rm sel},\qquad A^m=B^n.
\]
All displayed factors commute by the exact EX root/Cartier maps.
The base-preserving correction is consequently the literal map
\[
 C=A^3\mathcal D^{-3}
    =R_\lambda^{3n}\zeta_m^{3d}\widetilde F^{\,3n},
 \qquad
 C^{-1}=R_\lambda^{-3n}\zeta_m^{-3d}\widetilde F^{-3n}.
 \tag{DZ6.8}
\]
Its source translation is \(-3nR_{\rm div}\). On the cap
quotient, DZ6.7 therefore has endpoint \(C^{-1}\) in the
new-to-old direction and \(C\) in the old-to-new direction:
\(\mathcal D^3=A^3C^{-1}\). The scalar
\(R_\lambda^{3n}\) and the finite character are retained.

In the good branch \(R_{\rm div}\) lies in the identity formal
group. In the reduced branch it has component two, so
\(3R_{\rm div}\) lies in the identity formal group. The
central source map of \(\widetilde F^3\) is consequently the
identity in both cases, and its central line action is a nonzero
constant. The explicit root field and line map provide a formal
logarithm after this fixed central scalar is retained. This
logarithm extends through the native cap.

There is still an actual integral difference between this cap
logarithm and three annular logarithms. For the reduced branch,
with the reference Abel lifts retained,
\[
 3\log R_{\rm div}-\log(3R_{\rm div})
     =2\bigl(3\log G'-\log(3G')\bigr).
 \tag{DZ6.9}
\]
Indeed \(R_{\rm div}=2G'-[m]^{-1}((2m-1)G')\),
\(2m-1=3(2l+1)\), and the remaining terms lie in the
identity formal group, where their logarithms are additive.
For this formula take the literal sector lift
\(\log R_{\rm div}=2\log G'-m^{-1}\log((2m-1)G')\), using
the small formal logarithm in its second term. If the prescribed
annular lift differs by a period \(\Pi\), the exact simultaneous
replacement \(f\mapsto f+\Pi,\ h\mapsto h-ns\Pi\) preserves
the entire map EX.17. This carries that prescribed difference;
it does not change the actual annular marking.
The parenthesized quantity is the actual transported lattice
period fixed by this marking. Its period functions are retained;
it is not replaced by a constant complex number. For the complete
line logarithm its simultaneous physical-circle column is retained
by the same exact cochain replacement.

The ordinary root-annulus logarithm obeys
\[
 s_{\rm old}=s_{\rm new}+3\eta.
 \tag{DZ6.10}
\]
Thus the exact sector arrow EX.17 changes by the full
\(3\eta(nf+c_\lambda\delta)\) term, with the corresponding
reference-relative cochain. The following construction keeps
DZ6.9 and its line column in that square.

\par

\medskip\noindent\textbf{DZ6.4. The native full disc and the exact cap correction}

\par\nobreak

The native source addition has explicit regular charts. In the
exceptional chart \(X=T\alpha,Y=T\beta\),
\[
 \beta^2=T\alpha^3+R\alpha^2+D\alpha+1,\qquad
 \kappa=\frac{\beta-1}{\alpha}
        =\frac{D+R\alpha+T\alpha^2}{\beta+1}.
\]
The second expression is regular at \((\alpha,\beta)=(0,1)\);
the opposite chart uses the formal coordinate at \(O\).
The actual addition is
\[
 X'=\kappa^2-R-T\alpha,\qquad Y'=-T-\kappa X'.
\]
In the other source chart \(T=Xu,Y=Xv_0\), use
\[
 X=v_0^2-R-Du-u^2,\quad
 X'=u(D-2v_0+2u),\quad
 u'=\frac{X}{D-2v_0+2u},\quad v_0'=-u'-v_0+u.
 \tag{DZ6.11}
\]
At the two intersection points the denominator is a unit because
\((D-2v_0)(D+2v_0)=4s_0\ne0\). The inverse is the actual
opposite chord. These maps retain \(T\), \(w\) and the two
source resolution charts, throughout the native disc.

On the selected nodal fibre use
\[
 X=S^2-1,\qquad Y=S(S^2-w^2),\qquad G:\ S=-1.
\]
The full tangent-over-vertical factor becomes
\[
 f_2=\frac{(S+1)^2}{S+(1+w^2)/2}.
\]
It has value two on each branch \(S=\pm w\). Translation fixes
each of those node branches. Thus all six factors of \(B_{2,3}\)
give 64 there. For \(B_6=f_2^6 f_6(p-G)/f_6(p)\), the retained
Miller values on the node are
\(f_6=6+20w^2+6w^4\), which is 120 at \(w^2=3\);
the ratio is one. Hence, for both actual lines,
\[
                       (B_{\rm sel}/q)_{\rm node}
                            =64/(w^2-1)^3.
 \tag{DZ6.12}
\]
The actual Tate point \(P_G=(1+w)/(1-w)\), Tate-height unit
\(\kappa(0,w)=-(w^2-1)^2/(64w^6)\), and the two marked
Cartier factors
\[
 c_2=(1+w)^2/2,\qquad
 c_6=(1+w)^6/(6+20w^2+6w^4)
\]
are all nonzero analytic units on the chosen disc. Their full
Tate/Cartier maps therefore extend on every prescribed fan.
These statements use the actual parameters \(w\) near each
chosen square root of three, not an identification by fibre type.

The EX root, chord, invariant norm and full-reference annular
construction applies to DZ6.4 with \(b_0=Dz^l,c_0=Rz^{2l}\).
The certificate checks its universal identities before this
substitution. The selected source circle is disjoint from the
root and the I2 resolution. The fixed small scalar yields the
same compact contraction shells as EX.6. Thus the native model
is a proper smooth sphere family on the full disc, with its
literal identity closing.

Here is the cap map that implements the old/native comparison.
Let \(\mathcal G\) be the actual identity component of the
degree-zero line-translation group. Its exponential near the
identity is the exponential of the source dualizing vector and
physical \(\delta\) field. The regular EX vector and line map
extend this exponential through the resolved root. The map
\(\widetilde F^3\) has identity central source action and a
nonzero constant central line action. Retain the chosen central
line logarithm, divide by that central scalar to use the inverse
exponential near the identity, and then add it back. This gives
an extending logarithm \(\ell_3=\operatorname{Log}(\widetilde F^3)\).
It is fixed by the whole root deck, since its actual map and
central logarithm are fixed and the near-identity inverse is
unique. All these maps vary holomorphically on the native disc.

Set
\[
 \mathcal L_C=n\ell_3+3c_\lambda\delta,\qquad
 K_\eta=\operatorname{Exp}(\eta\mathcal L_C),\qquad
 \widehat J_\eta=J_\eta K_\eta.
 \tag{DZ6.13}
\]
The exponential is evaluated for the bounded complex times in
each fixed angular parameter chart; the source action is the
actual extending group flow. Since \(\operatorname{Exp}\mathcal
L_C=C\), the complete endpoint is
\[
 \widehat J_{\eta+1}
  =\mathcal D_{\rm old}^3 J_\eta C_{\rm new}K_\eta
  =A_{\rm old}^3\widehat J_\eta.
 \tag{DZ6.14}
\]
Its inverse has endpoint \(A_{\rm new}^{-3}\).
Consequently these are single-valued maps on the cap quotients.
They preserve the root function, the actual line, and physical
\(\delta\), including the full scalar in DZ6.8.

\par

\medskip\noindent\textbf{DZ6.5. Literal annular square and relative groupoid solve}

\par\nobreak

Choose one nonzero reference \(D_*\). Transport the entire old
ordinary-to-cap map at \(D_*\) by DZ6.5-\/-6, including its root
cut, character, source periods and every integral \(\delta\)
column. Construct the native reference map over the full native
disc from this transported map by the fixed-sector construction
EX.17-\/-19. The root norm branch uses the transported old value.
This specifies the whole native annular marking.

On a parameter angular chart and a base-annulus sector let
\(f=\operatorname{Log}\widetilde F\), and put
\[
                 \Pi_F=3f-\ell_3.
 \tag{DZ6.15}
\]
This is the actual period of the complete source/line Lie group,
because both exponentials are \(\widetilde F^3\).
Its source part in the reduced branch is the negative of DZ6.9;
its line part is the same retained period of the actual chord
lift. Pull the two annular cochains into the common native frame
and denote them by \(h_{\rm old}\), \(h_{\rm new}\).
Here this frame change includes DZ6.5c-\/-e: a tensor-reference
map is first conjugated by its complete ratio
\(f_P^a/(f_a f_{P^a})\). The new-to-old tensor arrow has the
literal \(+(a-1)\eta\delta\) contribution; its old-to-new
inverse \(I_D\) has the negative contribution. On adjacent
parameter charts these give the respective integral columns
\(+(a-1)j\delta\) and \(-(a-1)j\delta\). The transported
cochains use these directed frame arrows in full. The next
display is in the common point frame, after precisely these
conjugations.
Direct substitution of \(s_{\rm old}=s_{\rm new}+3\eta\)
and DZ6.13 in EX.17 gives
\[
 \widehat J_\eta^{-1}\Phi_{\rm old}
  =\operatorname{Exp}\!
     \left(\eta n\Pi_F+h_{\rm old}-h_{\rm new}\right)
      \Phi_{\rm new}I_D .
 \tag{DZ6.16}
\]
Here \(I_D\) is the actual old-to-new ordinary source/Cartier
map DZ6.5. This is a square of quotient maps with the same
ordinary input. The term \(3c_\lambda\delta\), including
both the scalar and character, cancels exactly between the two
displayed arrows. On a parameter-chart change
\(\eta_\beta=\eta_\alpha+j\), the indicated lifted translation
changes by \(jn\Pi_F\), together with the transported cochain
columns. This is its actual integral groupoid discrepancy.
It is retained rather than regarded as an extending root
logarithm.

For completeness, the relative solution uses the following
explicit sheaf and operators. On the full root atlas the
extending line-translation Lie bundle has
\[
 0\longrightarrow\mathcal O\delta\longrightarrow\mathcal V
       \longrightarrow\mathcal O(-1)\longrightarrow0.
 \tag{DZ6.17}
\]
Indeed the source dualizing vector pulls to \(z^k\) times the
root vector, with its exact character and \(0\le k<m\).
Taking its invariant holomorphic coefficients gives precisely
the original source Lie line. The invariant Cartier graph
arrows already proved identify its line extension; the physical
line is \(\mathcal O\delta\). Here are the actual splitting
and solution operators. Choose the source frames
\(e_\infty=T^{-1}e_0\) and write \(q_\infty=T^{-1}\),
retaining \(q=T-v\) for the selected node. On a fixed
overlap annulus expand the adjusted coefficient as
\[
 f(T,D)=\sum_{j\in\mathbb Z}f_j(D)T^j,\qquad
 f_j(D)=\frac1{2\pi i}\int_{|T|=\rho}
                     f(T,D)T^{-j-1}\,dT .
 \tag{DZ6.17c}
\]
For the source \(\mathcal O(-1)\) component take
\[
 u_0=-\sum_{j\ge0}f_jT^j,\qquad
 u_\infty=\sum_{j\le-1}f_j q_\infty^{-j-1}.
\]
Then \(T^{-1}u_\infty-u_0=f\) literally. For the
\(\mathcal O\delta\) component use the same \(u_0\) and
\(u_\infty=\sum_{j\le-1}f_jq_\infty^{-j}\); here
\(u_\infty-u_0=f\).
These series converge on the two chosen source-base charts
by their Laurent coefficient integrals.

The upper triangular transition of the actual Lie frames
is \(\widehat e_\infty=T^{-1}\widehat e_0+
a(T,D)\delta\). Its coefficient is the derivative of the
displayed Cartier arrow: in line-log coordinates it is the
source invariant vector applied to the logarithm of that
arrow\textquotesingle s full line multiplier, divided by \(2\pi i\).
Thus the tensor factor, root monomials and all Cartier units
are the functions being differentiated. Changing these
lifts to \(\widehat e_0+v_0\delta\) and
\(\widehat e_\infty+v_\infty\delta\) gives the exact equation
\[
              T v_\infty-v_0=-T a(T,D).
 \tag{DZ6.17d}
\]
For \(g=-Ta=\sum g_jT^j\), take
\(v_0=-\sum_{j\ge0}g_jT^j\) and
\(v_\infty=\sum_{j\le-1}g_j q_\infty^{1-j}\).
This is the explicit \(\mathcal O(1)\) Laurent solve, and it
splits DZ6.17 before the two component solves above.
All contours and Laurent projections are fixed, so these
operators preserve holomorphic parameter dependence and
a zero reference value.

Finally, after retaining the integral finite-group cochains,
let \(a_g\) be the remaining vector one-cocycle, with
\(a_{gh}=a_g+g a_h\). Its literal finite-group solution is
\[
                  v=-\frac1m\sum_h a_h,\qquad
                         g v-v=a_g .
 \tag{DZ6.17e}
\]
The equality follows by reindexing the entire sum using the
displayed cocycle identity. Averaging is applied only to these
complex vector cochains; every integral period and character
has been kept in the preceding adjustment.

Use the exponential sequence on this atlas and its entire finite
groupoid. At reduced central components take the actual
suspension power that moves the component to the identity before
the local logarithm. Every finite discrepancy retains
\(A^m=B^n\), its character and the line period in DZ6.15.
After the entire reference map has been matched at \(D_*\),
the difference of the integral base-direction two-cocycles is
zero there. It stays zero on each connected parameter chart:
all its values are continuous lattice-valued functions under
the actual source, root and Cartier identifications. On changing
parameter charts, DZ6.14 and DZ6.16 give exactly the integral
one-cochain \(jn\Pi_F\); its coboundary is the corresponding
change of representatives of those same two-cocycles.
Thus the zero class is transported around the whole annulus.
This records the complete mixed parameter/base discrepancy.

Subtract this integral cochain adjustment, apply the fixed
Laurent/Cauchy solution of DZ6.17 to the remaining vector
one-cocycle, and exponentiate. On each parameter chart this
gives an actual old/native family isomorphism. All its local
maps are the source addition/line maps above, the finite-root
maps, or the selected toric translations. At the cusp an
integral component reindexing is made first; the remaining
two characters are bounded holomorphic units. They act on
every prescribed Cox chart by its literal character map.
The inverses use the inverse group elements and the same
proper chart domains.

\par

\medskip\noindent\textbf{DZ6.6. Global physical-scalar factorization and closing}

\par\nobreak

The ratio of the local parameter-chart isomorphisms is a
global extending translation of the common marked family.
Here is its exact global group. Lift the marked zero section
through the nonzero-line covering away from the selected cusp.
Its source projection extends over every root; the regular-locus
criterion puts it on a smooth central component. The actual
suspension power moves it to the identity component, and the
degree-zero Hom-line evaluation gives its full line lift.
On generic fibres, constant factors of automorphy show that
this lift commutes with the actual suspension. Repeating the
finite-cap conjugacy \(m\) times and using strict contraction
forces the residual deck exponent to vanish. Thus all finite
characters and component choices have been retained.

The resulting source section is \(jG\) with \(j\in\mathbb Z\)
locally constant: these native surfaces have the same marked
E6 plus A1 lattice and primitive generator of height \(1/6\).
More explicitly, the actual simultaneous source resolution above
transports the unimodular rational-source intersection lattice
from a nearby \(D\ne0\) fibre. The E6 and A1 component lattice
has determinant six, and \(G\) meets the same nonidentity
components with corrections \(4/3\) and \(1/2\), while
\(G\cdot O=0\). Its height is \(2-4/3-1/2=1/6\).
Write \(\iota\) for its index in the free Mordell-\/-Weil lattice
and \(t_{\rm MW}\) for the torsion order. The integral
determinant formula gives
\[
 1=\frac{6\,h(G_{\rm prim})}{t_{\rm MW}^2},\qquad
 h(G)=\iota^2 h(G_{\rm prim})=\frac16,\qquad
                 \iota^2 t_{\rm MW}^2=1.
\]
Thus both positive integers are one. This identifies the
actual marked generator on the entire native disc.
The local covering lifts patch to this section: their differences
are integral suspension powers on the source base with its
selected point deleted. Its orbifold fundamental group is a
free product of the finite cap groups, and a homomorphism from
it to \(\mathbb Z\) is zero. The integer one-cocycle is
therefore removed by its actual deck powers, including at the
finite root.

The full valuation computation can be made directly in the
present Miller frames. Let \(U_G=B_{a,1}^{-1}\), the actual
meromorphic line lift covering \(+G\). Because \(B_{a,b}\)
is its literal inverse \(b\)-fold composite, DZ6.3 gives
\[
                 \varrho=B_{\rm sel}^{-1}/U_G^b=T^3/q.
 \tag{DZ6.17a}
\]
The ordinary I2 order of \(U_G\) is \(-3/b\).
Its finite-root order is \(2m/b\), by the full root coordinate
maps and the nonzero central coefficient already proved.
Its selected-cusp order is zero; its source Tate coordinate
there is a unit. There are no other fibre orders.
An endpoint lift has the form \(h(T)U_G^j\). Its invertibility
gives
\[
 \operatorname{ord}_0h=3j/b,\qquad
 \operatorname{ord}_\infty h=-2j/b.
\]
At the selected cusp the literal integer reindexing by \(B\)
gives \(\operatorname{ord}_{q=0}h=-k\): composing the
endpoint with \(B^k\) removes that line order. Source period
reindexing changes it by a unit only. The degree of this exact
rational divisor is zero, so \(j=bk\). All three orders of
\(h\) now equal those of \(\varrho^k\); their quotient is a
nonzero constant \(c\). The complete map is consequently
\[
              hU_G^{bk}=cB_{\rm sel}^{-k}.
 \tag{DZ6.17b}
\]
Since \(B=\lambda B_{\rm sel}\) is the actual quotient deck,
this acts by physical scalar \(c\lambda^k\).
Evaluation at an ordinary marked origin makes that scalar
holomorphic in the parameter. This computation includes the
finite integrality and character conditions before the global
equality, rather than erasing a finite central kernel.

Let \(F_+=cF_-\) be the actual two parameter-chart comparison
maps, with \(c\) this nonzero holomorphic scalar. On the two
simply connected overlap components choose \(\ell=\log c\),
and a fixed angular partition \(\chi_++\chi_-=1\). Put
\[
 b_+=\chi_-\ell,\qquad b_-=-\chi_+\ell,\qquad
 \omega=\bar\partial b_\pm,\qquad
 v(D)=\frac1\pi\int\frac{\omega(E)}{D-E}\,dA(E).
\]
A radial cutoff is one on the smaller annulus carrying the loop
and supported in the larger one. Then
\(g_\pm=\exp(b_\pm-v)\) are nonzero holomorphic functions,
\(g_+/g_-=c\), and
\[
                         g_+^{-1}F_+=g_-^{-1}F_- .
 \tag{DZ6.18}
\]
This is a single-valued actual family isomorphism over that
annulus. It includes all windings of \(c\) and the literal
physical action on the line and normal charts.

Both families are treated with their specified forward endpoint
maps \(E_{\rm old},E_{\rm new}\); for the native disc the latter
is its coordinate identity. Restrict the closed family map to
the cut parameter interval, calling it \(F_t\). Equation
DZ6.18 gives the actual square
\[
                 F_1E_{\rm old}=E_{\rm new}F_0 .
 \tag{DZ6.19}
\]
The physical return uses the inverse endpoint on both sides.
Transporting a native-disc connection through \(F_t\) therefore
conjugates the complete oriented old return to that of the
native disc. The latter extends over the disc and is zero.
This proves the zero return for the specified old family with
the whole transported reference marking at either \(w^2=3\).
The exact original-moduli coordinate at infinity is
\[
                 \epsilon=\frac1r
                    =\frac{D^2}{3-2w^2}.
 \tag{DZ6.20}
\]
Since \(D\) is a simple local coordinate at either point and
the denominator has value \(-3\), a positive \(w\)-loop has
degree two over the positive original infinity meridian.
The constructed family map concerns this actual marked
\(w\)-loop, with the source and cap arrows already specified.
No shorter original-parameter endpoint or return is assigned.
The full source, line, cap and endpoint maps above prove the stated conclusion with every prescribed annular integer retained.

\par

\medskip\noindent\textbf{Executed exact certificate}

\par\nobreak

The exact program \path{code/part3/e6_d_zero_disc.py} and output \path{code/part3/e6_d_zero_disc.json} execute 521 identities and finite signed instances. They include the source and nodal Cartier charts, both root and tensor weights, complete cap endpoint, retained period, all three Laurent operators, finite vector averaging, global valuation equation and exact degree-two parameter. The complete geometric family proof is above. The preceding coordinate-only and frame-clarification stages remain under \path{records/attempts/e6_d_zero_disc/}.

\end{proof}
